\chapter{PCR (Polymerase Chain Reaction)}

\section*{What Is This Test?}
Polymerase chain reaction (PCR) is a revolutionary molecular biology technique that exponentially amplifies specific DNA or RNA sequences to detectable levels, enabling the detection and quantification of infectious pathogens, genetic mutations, and human gene sequences from minute quantities of starting material. In clinical infectious disease diagnostics, PCR (and reverse-transcription PCR/RT-PCR for RNA targets, collectively referred to as nucleic acid amplification tests/NAATs) allows the direct detection of pathogen nucleic acid in clinical specimens without the need for culture. Applications include: detection and quantification of viral loads (HIV, HBV, HCV, CMV, EBV, SARS-CoV-2), diagnosis of bacterial infections that are difficult or slow to culture (\textit{Chlamydia trachomatis}, \textit{Neisseria gonorrhoeae}, \textit{Mycobacterium tuberculosis}, \textit{Bordetella pertussis}, \textit{Mycoplasma pneumoniae}), detection of antimicrobial resistance genes (\textit{mecA} for MRSA, \textit{vanA/vanB} for VRE, \textit{bla}$_{KPC}$/\textit{bla}$_{NDM}$ for carbapenemases), rapid diagnosis of meningitis/encephalitis (multiplex FilmArray ME Panel), respiratory pathogen panels, gastrointestinal pathogen panels, and blood culture identification panels. PCR has transformed clinical microbiology by providing results in hours rather than days or weeks, improving patient outcomes through earlier targeted therapy and infection control interventions \cite{yang2004pcrbased}.

\section*{Alternative Names}
\begin{itemize}
  \item NAAT (Nucleic Acid Amplification Test)
  \item PCR Test
  \item Molecular Testing
  \item DNA Amplification
  \item RT-PCR
  \item qPCR
  \item Multiplex PCR
  \item Real-Time PCR
  \item Rapid Molecular Testing
\end{itemize}
\section*{Principle}
PCR involves cyclic temperature changes that drive repeated rounds of DNA synthesis: (1) Denaturation: The double-stranded DNA template is heated to 94--98\textdegree{}C, separating the strands. (2) Annealing: The temperature is lowered to 50--65\textdegree{}C, allowing short synthetic oligonucleotide primers (18--25 base pairs) to hybridize (anneal) to complementary sequences flanking the target region on each single-stranded DNA template. (3) Extension: The temperature is raised to 72\textdegree{}C (the optimum for \textit{Taq} DNA polymerase, a thermostable enzyme isolated from \textit{Thermus aquaticus}), and the polymerase synthesizes a new DNA strand complementary to the template by extending from the 3' end of the primer, incorporating deoxynucleotide triphosphates (dNTPs) \cite{saiki1988primer}. Each thermal cycle doubles the number of target DNA copies; after 30--40 cycles, a single target molecule generates >10$^9$ copies, sufficient for detection. Real-time PCR (quantitative PCR/qPCR): Incorporates a fluorescent reporter molecule (TaqMan probe, molecular beacon, or SYBR Green intercalating dye) that generates a fluorescent signal proportional to the amount of amplified product during each cycle in real time. The cycle threshold (Ct) is the cycle number at which the fluorescent signal exceeds the background threshold; the Ct is inversely proportional to the amount of target nucleic acid in the starting specimen (lower Ct = higher starting concentration) \cite{espy2006realtime}. Quantitative results are calibrated using external standards of known concentration and reported as copies/mL, IU/mL (WHO International Standards for HIV, HBV, HCV, CMV, EBV), or genome equivalents/mL \cite{bustin2009miqe}. Reverse-transcription PCR (RT-PCR) is used for RNA targets (HIV, HCV, SARS-CoV-2, influenza, RSV): viral RNA is first reverse-transcribed into complementary DNA (cDNA) by the enzyme reverse transcriptase, and the cDNA is then amplified by standard PCR. Nested PCR: A second round of amplification using internal primers nested within the first-round PCR product increases sensitivity and specificity. Multiplex PCR: Multiple primer pairs are included in a single reaction to simultaneously amplify and detect multiple pathogen targets (e.g., BioFire FilmArray multiplex panels). Digital PCR: Absolute quantitation without standard curves by partitioning the sample into thousands of individual nanoliter-scale reactions and counting the number of positive partitions.

\section*{History}
PCR was invented by Kary Mullis in 1983 (Nobel Prize 1993), who conceived the idea during a drive along the California coast. The first clinical PCR test (HIV proviral DNA) was introduced in the late 1980s. Real-time PCR (TaqMan) was invented in the 1990s by researchers at Roche. FDA-cleared multiplex respiratory panels became available in the 2010s and represent the current state-of-the-art syndromic molecular testing. The COVID-19 pandemic (2020--2023) drove unprecedented global adoption of RT-PCR testing, with billions of SARS-CoV-2 tests performed. The development of fully automated, cartridge-based, sample-to-answer molecular platforms (Cepheid GeneXpert, Roche cobas Liat, BioFire FilmArray, Abbott ID NOW) brought PCR into point-of-care settings and small laboratories without molecular expertise.

\section*{Why This Test Is Ordered}
PCR is ordered for hundreds of specific clinical indications detailed in the respective organ-system and pathogen-specific chapters. General categories include:
\begin{itemize}
  \item Diagnosis of viral infections: HIV, HCV, HBV, CMV, EBV, HSV, VZV, influenza, RSV, SARS-CoV-2, adenovirus, enterovirus, norovirus, rotavirus, HPV (cervical cancer screening), etc.
  \item Diagnosis of bacterial infections where culture is slow, insensitive, or not routinely performed: \textit{Chlamydia trachomatis}, \textit{Neisseria gonorrhoeae}, \textit{Mycobacterium tuberculosis} (Xpert MTB/RIF), \textit{Bordetella pertussis}, \textit{Mycoplasma pneumoniae}, \textit{Chlamydia pneumoniae}, \textit{Legionella pneumophila}
  \item Quantitation of viral loads for monitoring of chronic viral infections and response to therapy: HIV RNA, HBV DNA, HCV RNA, CMV DNA, EBV DNA, BK virus (in renal transplant)
  \item Detection of antimicrobial resistance genes: \textit{mecA} (MRSA nares screen), \textit{vanA/vanB} (VRE rectal screen), carbapenemase genes (\textit{bla}$_{KPC}$, \textit{bla}$_{NDM}$, \textit{bla}$_{VIM}$, \textit{bla}$_{IMP}$, \textit{bla}$_{OXA-48-like}$), ESBL genes
  \item Multiplex syndromic panels for rapid diagnosis of sepsis (blood culture identification panels: BioFire BCID2, Verigene BC-GP/BC-GN, ePlex BCID), meningitis/encephalitis (FilmArray ME), respiratory infections (FilmArray RP2.1, Luminex NxTAG RPP), gastrointestinal infections (FilmArray GI, Luminex xTAG GPP), and joint infections (BioFire Joint Infection Panel)
  \item Confirmatory testing: PCR for \textit{T. pallidum} (syphilis) from lesion exudate, \textit{Toxoplasma gondii} in immunocompromised/neonates, and molecular typing of pathogens for epidemiologic investigations
\end{itemize}

\section*{Primary vs Secondary Test}
PCR has become the primary diagnostic method for many infections (chlamydia, gonorrhea, influenza, SARS-CoV-2, HIV viral load monitoring, CMV/EBV in transplant) \cite{caliendo2013better}. For others (TB, pertussis, HSV encephalitis, \textit{Clostridioides difficile}), PCR is a primary rapid test complementary to culture. In still others (bacterial wound infections, pneumonia, urinary tract infections), PCR is a secondary adjunct to culture.

\section*{How This Test Is Performed}
Specimen requirements are assay-specific and detailed in the relevant test chapters. General principles: (1) Specimens for PCR should be collected as early in the illness as possible when pathogen load is highest. (2) Swabs for PCR should be placed in the manufacturer's specific transport medium (NOT routine Amies transport medium for bacteria, which may contain PCR inhibitors). (3) Blood for viral load quantitation should be collected in EDTA (purple-top; for plasma) or PPT (white-top; for plasma separation); heparin inhibits PCR. (4) Tissue and body fluid specimens should be placed in sterile containers without formalin (formalin cross-links nucleic acids and prevents amplification). (5) Freezing at -70\textdegree{}C is the preferred storage method for specimens that will undergo PCR, though some assays tolerate room temperature or refrigeration for limited periods.

\section*{What Preparation Is Needed}
No special patient preparation. Pre-analytical variables including recent antibiotic/antiviral therapy should be documented, as they may affect nucleic acid detection but generally do not produce false-negative PCR results as quickly as they do for culture. Patients should avoid eating, drinking, or using mouthwash 30 minutes before oropharyngeal or nasopharyngeal specimens. Heparin anticoagulants should be avoided for blood specimens.

\section*{Contraindications and Precautions}
\begin{itemize}
  \item PCR is exquisitely sensitive to contamination: even a single molecule of amplified product (amplicon) from a previous reaction, or target DNA introduced from the environment, can cause false-positive results. Clinical laboratories employ strict unidirectional workflow (reagent preparation $\rightarrow$ specimen processing $\rightarrow$ amplification/detection), separate dedicated areas, and UV decontamination. False-negative PCR results may occur because of: (a) PCR inhibitors in the specimen (heparin, hemoglobin, lactoferrin, IgG, urea, formalin); (b) low organism load below the assay's limit of detection; (c) mutations in the primer or probe binding sites (important for rapidly mutating RNA viruses like HIV, HCV, SARS-CoV-2); (d) improper specimen collection, transport, or storage leading to nucleic acid degradation; (e) the presence of the pathogen at a site not sampled (e.g., CSF PCR negative for HSV in early HSV encephalitis). PCR detects nucleic acid from both viable and nonviable organisms; therefore, PCR may remain positive for days to weeks after successful treatment (<3 weeks for most bacteria; for EBV/CMV DNA, low-level positivity may persist for months). This limitation has implications for test of cure. CT values from real-time PCR should be interpreted with caution: they are not standardized across assays and do not necessarily correlate perfectly with clinical severity, though very low Ct values (e.g., <20) generally indicate high organism burden. Quantitation of viral loads has been standardized with WHO International Standards for HIV, HBV, HCV, CMV, and EBV. Not all PCR assays are quantitative; many are qualitative (detected/not detected). A positive qualitative PCR does not provide a measure of organism burden and cannot be used to monitor response to therapy.
\end{itemize}
\section*{Risks}
No risks from the PCR test itself. Risks are associated with specimen collection (venipuncture, lumbar puncture, swab, biopsy). The major clinical risk of RT-PCR is false reassurance from a false-negative result, particularly in early infection when pathogen load is low (the window period).

\section*{What Happens After the Test}
Turnaround times range from 15--90 minutes (rapid molecular platforms: GeneXpert, FilmArray, cobas Liat) to 2--6 hours (standard real-time PCR) to 24--48 hours (send-out reference laboratory PCR assays). Results are reported as: "Detected" or "Not Detected" for qualitative PCR; a quantitative viral load in copies/mL or IU/mL for qPCR; and cycle threshold (Ct) values may or may not be reported.

\section*{Understanding Results}

\subsection*{Below Normal}
\begin{itemize}
  \item \textbf{Not detected (or Target not detected, TND):} No pathogen-specific nucleic acid detected. Interpretation depends on the clinical context, timing of specimen collection (window period), specimen adequacy, and assay sensitivity. A negative PCR does NOT categorically exclude infection. For quantitative assays, "not detected" is the goal of therapy for HIV, HCV, and HBV DNA suppression
\end{itemize}

\subsection*{Above Normal}
\begin{itemize}
  \item \textbf{Detected (positive):} Pathogen nucleic acid detected. For qualitative assays, indicates the presence of the target organism/mutation. For quantitative assays, reported as the viral load in the appropriate units. Must be interpreted in the clinical context (distinguishing active infection from colonization, latency, or post-treatment nucleic acid persistence)
  \item \textbf{Quantitative result (viral load):} For HIV, HBV, HCV, CMV, EBV, BK virus. The numeric result (in copies/mL or IU/mL) is compared to prior results (trend) and to clinical thresholds for treatment decisions
  \item \textbf{Ct value (cycle threshold):} Semi-quantitative inversely related to organism burden. Ct <20 = high load; Ct 20--30 = moderate; Ct 30--40 = low load. Ct values are assay-specific and should not be directly compared between different assays
\end{itemize}

\section*{Normal Results}
\textbf{Not detected} or \textbf{Target not detected} (qualitative). For quantitative assays in uninfected individuals: \textbf{Not detected} or \textbf{<LLOQ}. For HIV-infected patients on suppressive ART: \textbf{Not detected} (<20--50 copies/mL depending on assay).

\section*{Follow-up Tests}
Follow-up depends on the specific pathogen and clinical scenario, as detailed in the individual test chapters. General approach:
\begin{itemize}
  \item \textbf{Culture with antimicrobial susceptibility:} When PCR detects a pathogen for which phenotypic AST is required (bacteria, \textit{Candida}). PCR identifies the organism; culture provides the isolate for AST
  \item \textbf{Genotypic resistance testing:} For HIV (protease, RT, integrase sequencing), HCV (NS3, NS5A, NS5B), CMV (UL97, UL54), HSV/VZV (thymidine kinase, DNA polymerase), TB (GeneXpert, line probe assay), and \textit{Staphylococcus} (\textit{mecA}, \textit{vanA/vanB})
  \item \textbf{Repeat quantitative PCR:} For monitoring of viral loads at defined intervals (HIV: every 3--6 months; CMV post-transplant: weekly; HCV: baseline, week 4 on treatment, end of treatment, SVR12)
  \item \textbf{Confirmatory testing with alternative method:} When the initial PCR result is unexpected, discordant with the clinical picture, or has significant therapeutic implications
\end{itemize}

\section*{References}
\bibliographystyle{unsrtnat}
\bibliography{references}

\clearpage
\section*{Regulatory Status and Authorization Date}
This chapter describes a test category, not a specific commercial product. FDA, CDSCO, EU IVDR, and MHRA approval or registration dates therefore cannot be assigned without the assay manufacturer, product name, instrument, and intended use. For laboratory-developed tests, an individual agency approval date may not exist. See the Preface for the interpretation of regulatory status.

\clearpage
